\documentclass[11pt]{article}

% ============================================================
%  Packages
% ============================================================
\usepackage[margin=0.85in,top=0.75in,bottom=0.85in]{geometry}
\usepackage{amsmath}
\usepackage{amssymb}
\usepackage{booktabs}
\usepackage{array}
\usepackage[skip=4pt]{caption}
\usepackage{enumitem}
\usepackage{titlesec}
\usepackage{xcolor}
\usepackage[most]{tcolorbox}
\usepackage{ragged2e}
\usepackage[hidelinks]{hyperref}

% ============================================================
%  Color tokens
% ============================================================
\definecolor{moss}{HTML}{758C58}
\definecolor{mossdk}{HTML}{55663F}   % darker moss for headings
\definecolor{ink}{HTML}{1C1C1E}      % near-black body text
\definecolor{rulegray}{HTML}{D8D8D2} % light gray hairlines
\definecolor{mutegray}{HTML}{6E6E73} % captions / secondary text
\definecolor{mossbg}{HTML}{F2F4EE}   % faint moss wash (moss!5-ish)

\color{ink}

% ============================================================
%  Typographic setup
% ============================================================
\setlength{\parindent}{0pt}
\setlength{\parskip}{0.38em}
\setlist{itemsep=2pt, topsep=3pt, parsep=0pt}

% Section headings: moss, bold, with a hairline underneath
\titleformat{\section}
  {\Large\bfseries\color{mossdk}}
  {}{0pt}
  {}
  [{\vspace{2pt}\color{rulegray}\titlerule[0.8pt]}]
\titlespacing*{\section}{0pt}{11pt}{6pt}

\titleformat{\subsection}
  {\normalsize\bfseries\color{ink}}
  {}{0pt}{}
\titlespacing*{\subsection}{0pt}{9pt}{2pt}

% Kill default section numbers globally (we head manually)
\setcounter{secnumdepth}{0}

% ============================================================
%  Hero-stat callout: big moss number + small caption
%  #1 = big number   #2 = caption
% ============================================================
\newcommand{\herostat}[2]{%
  \begin{minipage}[t]{0.315\linewidth}%
    \centering
    {\fontsize{30}{32}\selectfont\bfseries\color{moss} #1}\\[5pt]
    {\footnotesize\color{mutegray}\setlength{\baselineskip}{11pt} #2\par}%
  \end{minipage}%
}

% Row wrapper for three hero stats, framed top and bottom by hairlines
\newenvironment{herorow}
  {\par\vspace{4pt}\noindent\color{rulegray}\rule{\linewidth}{0.8pt}\par
   \vspace{9pt}\noindent}
  {\par\vspace{9pt}\noindent\color{rulegray}\rule{\linewidth}{0.8pt}\par\vspace{4pt}}

% ============================================================
%  Result "card": moss frame, faint moss fill, hard corners
% ============================================================
\newtcolorbox{resultcard}[1][]{
  colframe=moss,
  colback=mossbg,
  boxrule=0.5pt,
  sharp corners,
  left=10pt, right=10pt, top=8pt, bottom=8pt,
  fonttitle=\bfseries\color{mossdk},
  coltitle=mossdk,
  title={#1},
  before skip=8pt, after skip=8pt,
}

% Synthesis card: inverted emphasis, moss title bar
\newtcolorbox{synthcard}[1][]{
  enhanced,
  colframe=mossdk,
  colback=white,
  boxrule=0.9pt,
  sharp corners,
  left=12pt, right=12pt, top=10pt, bottom=10pt,
  title={#1},
  fonttitle=\bfseries\large\color{white},
  coltitle=white,
  colbacktitle=mossdk,
  attach boxed title to top left={xshift=0pt, yshift=0pt},
  boxed title style={sharp corners, boxrule=0pt},
  before skip=10pt, after skip=10pt,
}

% Small moss tag label used inline
\newcommand{\mosstag}[1]{{\footnotesize\bfseries\color{moss}\MakeUppercase{#1}}}

\begin{document}

% ============================================================
%  Title block
% ============================================================
\begin{flushleft}
{\color{rulegray}\rule{\linewidth}{1.2pt}}\\[10pt]
{\LARGE\bfseries\color{ink} Reading the Prompt, Not the Economics}\\[6pt]
{\large\color{mossdk} Structural Estimation of an AI Model's Risk Preferences}\\[3pt]
{\normalsize\color{mutegray} A progress report on Phases 1 and 2, with the Phase 3 plan}\\[11pt]
{\color{rulegray}\rule{\linewidth}{0.8pt}}\\[7pt]
{\small\color{ink}\textbf{J. Blakeslee}, Pardee RAND Graduate School \hfill
 AI Sprint Practicum, Course 1012 \quad$\cdot$\quad July 2026}\\[3pt]
{\color{rulegray}\rule{\linewidth}{1.2pt}}
\end{flushleft}

\vspace{6pt}

% One-line thesis, featured
\begin{tcolorbox}[
  colframe=moss, colback=mossbg, boxrule=0.5pt, sharp corners,
  left=12pt, right=12pt, top=9pt, bottom=9pt, before skip=2pt, after skip=10pt]
{\large\itshape\color{ink} Across two phases, Claude's risky choices track the surface features
of a prompt, the probability figure and the word ``loss,'' more than the underlying economics.}
\end{tcolorbox}

This project treats a large language model as an economic agent and recovers its preference
primitives from its choices. Two phases are complete. Each elicits roughly a thousand to
two-and-a-half thousand controlled binary lottery choices from a pinned model snapshot, then fits a
structural discrete-choice model to the result. Phase 1 asks how the model weighs risk. Phase 2 asks
whether the wording of a choice, holding the money fixed, moves the answer. Both return sharp,
quantified, human-comparable results. Phase 3, planned below, asks where the biases come from.

% ============================================================
%  PHASE 1
% ============================================================
\section{Phase 1: How Claude Weighs Risk}

{\small\color{mutegray} \texttt{claude-haiku-4-5-20251001} \;$\cdot$\; 1{,}200 controlled choices
\;$\cdot$\; forced binary choice between a sure amount and a two-outcome gamble}

\begin{herorow}
\herostat{2.69}{probability-weighting curve; humans about 0.7 and bent the opposite way}%
\hfill
\herostat{0.59}{risk-aversion once weighting is controlled; squarely human-plausible}%
\hfill
\herostat{22.8 vs.\ 5.0}{how strongly the win-probability drives choice versus the payoff size}
\end{herorow}

A first, simple risk model fit badly. A textbook constant-relative-risk-aversion (CRRA) expected
utility model implied a risk-aversion of 0.19 and could not reproduce the choices. The diagnostic
showed why. As the win probability rose from low to high, gamble-taking climbed from 8\% to 86\%,
while the gambles' actual expected value stayed roughly flat. Claude was pricing the probability and
nearly ignoring the dollar amounts.

An extended model, adding a probability-weighting curve and a position term, fit far better: a
likelihood-ratio chi-square of 617 on 2 degrees of freedom. Its probability-weighting curve is an
extreme S-shape at 2.69, the opposite of the human inverse-S. Claude underweights long shots instead
of overweighting them, so it will refuse a long shot at almost any prize. Once that weighting is
controlled, its risk-aversion of 0.59 is ordinary and human-plausible; the earlier 0.19 was an
artifact of the wrong model. Choice noise roughly halved, from 0.26 to 0.10, once the model was right.

\begin{resultcard}[Key finding]
Claude prices the probability of winning and nearly ignores the size of the prize. Its
probability-weighting curve is an extreme S-shape pointed opposite to the human one: it underweights
long shots rather than overweighting them, and will pass on a high-value long shot at almost any
prize.
\end{resultcard}

\begin{table}[h]
\centering
\small
\caption*{\small\textbf{M2 parameter estimates.} CRRA value, Prelec probability weighting, and a
position term. $n = 1{,}200$.}
\begin{tabular}{@{}l l l@{}}
\toprule
\textbf{Parameter} & \textbf{Estimate (se)} & \textbf{Interpretation} \\
\midrule
$r$ (curvature)       & 0.592 (0.028) & Risk aversion, human-plausible range \\
$\gamma$ (Prelec)     & 2.693 (0.166) & Extreme S-shape; long shots underweighted \\
$\mu$ (Fechner noise) & 0.104 (0.008) & Choice noise, less than half of the baseline \\
$\delta$ (position)   & 1.257 (0.160) & About $3.5\times$ odds pull toward the option listed second \\
\bottomrule
\end{tabular}
\end{table}

% ============================================================
%  PHASE 2
% ============================================================
\section{Phase 2: Framing and the Fear of Losses}

{\small\color{mutegray} \texttt{claude-haiku-4-5-20251001} \;$\cdot$\; 2{,}480 controlled choices
\;$\cdot$\; one terminal-wealth lottery, shown in three logically identical frames}

\begin{herorow}
\herostat{40\% $\rightarrow$ 10\%}{gamble-taking on the same lottery when the downside is worded as a
loss instead of getting nothing}%
\hfill
\herostat{3.23}{estimated fear-of-losses, above the classic human figure of about 2.25}%
\hfill
\herostat{100\%}{a built-in control, where one option was plainly larger, answered correctly every time}
\end{herorow}

The same terminal-wealth lottery was shown three ways that a rational agent must treat identically: as
pure gains, in neutral wording, and as a mix where the downside is called a loss. Gamble-taking was
40\% in the gain wording, 30\% neutral, and 10\% when the downside was a loss, a gap of about 0.30
with a p-value on the order of $10^{-47}$. Holding each individual lottery fixed, 88\% of them show the
effect, so it is not composition. This violates frame invariance: identical money, different words,
different choice.

The structural estimate puts the fear-of-losses at 3.23, above the canonical human figure of about
2.25. The dominant-choice control, a certain amount versus a strictly larger certain amount, was
answered correctly 100\% of the time. Claude reads dollar magnitudes perfectly well when no
probability is involved; the trouble is specific to choices under risk, and here, specific to the word
``loss.''

\begin{resultcard}[Key finding]
Two lotteries with byte-for-byte identical terminal wealth draw 40\% versus 10\% gamble-taking based
only on whether the downside is called a loss. The estimated loss-aversion coefficient, 3.23, exceeds
the canonical human 2.25. Wording moves the choice; the money does not change.
\end{resultcard}

\begin{table}[h]
\centering
\small
\caption*{\small\textbf{Structural estimates under the stated-reference assumption} ($\phi = 1$,
$n = 2{,}480$).}
\begin{tabular}{@{}l l l@{}}
\toprule
\textbf{Parameter} & \textbf{Estimate (se)} & \textbf{Interpretation} \\
\midrule
$\lambda$ (loss aversion) & 3.234 (0.171)  & Losses loom about $3.2\times$ larger than equal gains \\
$\alpha$ (curvature)      & 0.503 (0.024)  & Diminishing sensitivity, human-plausible \\
$\gamma$ (Prelec)         & 1.516 (0.059)  & S-shape; long shots underweighted \\
$\mu$ (Fechner noise)     & 0.131 (0.009)  & Choice noise \\
$\delta$ (position)       & $-0.363$ (0.104) & Pull away from the gamble when listed second \\
\bottomrule
\end{tabular}
\end{table}

\subsection{An honest identification limit}
The data cannot fully separate ``the reference point moved'' from ``an extreme fear of losses.'' The
mixed-frame suppression admits two nearly observationally equivalent readings, so the clean headline is
the assumption-free frame gap, and the fear-of-losses figure of 3.23 is reported under the standard
stated-reference assumption. One further caution: the position term flipped sign between phases, from
$+1.26$ in Phase 1 to $-0.36$ here, so position bias is a design artifact rather than a stable trait
of the model.

% ============================================================
%  SYNTHESIS
% ============================================================
\section{Synthesis: Two Phases, One Thread}

\begin{synthcard}[The emerging portrait]
Claude's choices under risk are governed by salient surface features of the prompt, the probability
number in Phase 1 and the loss wording in Phase 2, more than by an integrated computation over final
wealth. It reads magnitudes fine when nothing distracts it, as the 100\% control shows. So this is
selective attention, not an inability to do arithmetic.
\end{synthcard}

% ============================================================
%  PHASE 3
% ============================================================
\section{Phase 3: Where Do the Biases Come From?}

This is the payoff the whole project was built toward. Phases 1 and 2 measured the biases. Phase 3
asks the attribution question: are these distortions written into the model during pretraining, or
installed by the final alignment step?

\textbf{Method.} Run the identical Phase 1 and Phase 2 instruments on matched pairs of open models,
one \emph{before} its final training step and one \emph{after}, that is, base versus instruct. Read
each choice probability directly from the model's own token probabilities rather than from sampled
text. If a bias appears only in the after-training model, the training instilled it. A format-control
arm, running the after-training model in both prompt formats, rules out a wording confound.

\textbf{Models.} Qwen2.5-7B goes first, then Llama-3.1-8B, with OLMo-2 as a stretch. OLMo-2 publishes
its staged checkpoints, base, then supervised, then preference-tuned, which would let the bias be
pinned to a \emph{specific} training stage.

\begin{resultcard}[The question Phase 3 answers]
Is a model's economic character learned in pretraining, or written in by the alignment step at the
end?
\end{resultcard}

\textbf{Status and practicalities.} The pipeline is built and already validated on simulated data,
where it recovers known values near-perfectly. It runs on a free Colab GPU, because no hosted service
serves the needed before-training models with token probabilities.

\vspace{6pt}
{\color{rulegray}\rule{\linewidth}{0.8pt}}\\[3pt]
{\footnotesize\color{mutegray} Estimates are properties of the pinned snapshots under this
unincentivized forced-choice protocol; they are not claimed as universal traits of language models.}

\end{document}
